% FAC-EN-U0024
% French source: fac_body.tex lines 873--892.
% Authority: PDF physical pp. 23--24; printed pp. 218--219 (no. 26 complete).
% Printed-source repair: FAC-FR-C0010.

\medskip
\phantomsection\label{fac:I:3:no:26}
\textbf{26. Cohomology of a closed subspace.}
Let $\mathcal{F}$ be a sheaf on the space $X$, and let $Y$ be a subspace of $X$. Let $\mathcal{F}(Y)$ be the sheaf induced by $\mathcal{F}$ on $Y$, in the sense of \hyperref[fac:I:1:no:4]{no.~4}. If
$\mathfrak{U}=\{U_i\}_{i\in I}$ is a covering of $X$, the sets
$U'_i=Y\cap U_i$ form a covering $\mathfrak{U}'$ of $Y$. If
$f_{i_0\ldots i_q}$ is a section of $\mathcal{F}$ over $U_{i_0\ldots i_q}$, its restriction to
\[
U'_{i_0\ldots i_q}=Y\cap U_{i_0\ldots i_q}
\]
is a section of $\mathcal{F}(Y)$.\footnote{The printed source places a stray comma between $U'_{i_0\ldots i_q}$ and the defining equality. It is removed here and in the corrected French layer.}
Restriction is a homomorphism
\[
\rho:C(\mathfrak{U},\mathcal{F})
\longrightarrow C(\mathfrak{U}',\mathcal{F}(Y))
\]
which commutes with $d$, and hence defines
\[
\rho^*:H^q(\mathfrak{U},\mathcal{F})
\longrightarrow H^q(\mathfrak{U}',\mathcal{F}(Y)).
\]
If $\mathfrak{U}<\mathfrak{V}$, then $\mathfrak{U}'<\mathfrak{V}'$, and
\[
\rho^*\circ\sigma(\mathfrak{U},\mathfrak{V})
=\sigma(\mathfrak{U}',\mathfrak{V}')\circ\rho^*.
\]
Thus, by passage to the limit over $\mathfrak{U}$, the homomorphisms $\rho^*$ define homomorphisms
\[
\rho^*:H^q(X,\mathcal{F})
\longrightarrow H^q(Y,\mathcal{F}(Y)).
\]

\phantomsection\label{fac:I:3:no:26:prop:8}
\textsc{Proposition 8.} \textit{Suppose that $Y$ is closed in $X$ and that $\mathcal{F}$ is zero outside $Y$. Then the homomorphism}
\[
\rho^*:H^q(X,\mathcal{F})\longrightarrow H^q(Y,\mathcal{F}(Y))
\]
\textit{is bijective for every $q\geq0$.}

The \hyperref[fac:I:3:no:26:prop:8]{proposition} follows from the following two facts:

\begin{enumerate}[label=(\alph*),leftmargin=2.5em]
\item Every covering $\mathfrak{W}=\{W_i\}_{i\in I}$ of $Y$ is of the form $\mathfrak{U}'$, where $\mathfrak{U}$ is a covering of $X$.

Indeed, it is enough to set $U_i=W_i\cup(X-Y)$, since $Y$ is closed in $X$.

\item For every covering $\mathfrak{U}$ of $X$, the homomorphism
\[
\rho:C(\mathfrak{U},\mathcal{F})
\longrightarrow C(\mathfrak{U}',\mathcal{F}(Y))
\]
is bijective.

Indeed, this follows from \hyperref[fac:I:1:no:5:prop:5]{Proposition~5} of \hyperref[fac:I:1:no:5]{no.~5}, applied to $U_{i_0\ldots i_q}$ and the sheaf $\mathcal{F}$.
\end{enumerate}

\hyperref[fac:I:3:no:26:prop:8]{Proposition~8} may also be expressed as follows. If $\mathcal{G}$ is a sheaf on $Y$, and if $\mathcal{G}^X$ is the sheaf obtained by extending $\mathcal{G}$ by zero outside $Y$, then
\[
H^q(Y,\mathcal{G})=H^q(X,\mathcal{G}^X)
\qquad(q\geq0).
\]
In other words, the identification of $\mathcal{G}$ with $\mathcal{G}^X$ is compatible with passage to cohomology groups.
